\reading{MR-16}{The Vasicek and Gauss+ Models}{defer}{2}
% ==========================================================

\begin{mrkeybox}[title={Learning objectives in this reading}]
\small
\begin{enumerate}[label=\textbf{\alph*.},leftmargin=1.7em]
\item Describe the structure of the Gauss+ model and discuss the implications of this structure
      for the model's ability to replicate empirically observed interest rate dynamics.
\item Compare and contrast the dynamics, features, and applications of the Vasicek model and the
      Gauss+ model.
\item Calculate changes in the short-term, medium-term, and long-term interest rate factors
      under the Gauss+ model.
\item Explain how the parameters of the Gauss+ model can be estimated from empirical data.
\end{enumerate}
\end{mrkeybox}

\section{Basics}

\term{Interest rate modeling} is crucial for risk management and for pricing interest rate
securities and derivatives. Good models explain the term structure and are consistent with
implied volatilities.

\begin{mrkeybox}[title={Objectives of a good model}]
\begin{itemize}
\item \term{Consistency with market prices} --- must match bond and derivative prices.
\item \term{Market expectations} --- should reflect mean reversion and volatility trend.
\end{itemize}
\end{mrkeybox}

\subsection{The Vasicek model, recalled}
The Vasicek model uses a \term{single factor} to model how short-term interest rates change over
time. Its main characteristic is mean reversion, where the short-term rate $r_t$ tends to revert
to its long-term average $\theta$. It is useful for simple applications like pricing zero-coupon
bonds and applying basic hedging techniques.

\begin{mrtrapbox}[title={Its drawbacks}]
\begin{enumerate}
\item It fails to replicate complex term structure shapes and real-world volatility patterns.
\item It fails to account for observed changes in volatility across maturities.
\item It fails to adequately incorporate macroeconomic trends and monetary policy shifts.
\end{enumerate}
\end{mrtrapbox}

\subsection{How interest rates are observed in the market}

\begin{center}
\begin{tikzpicture}[every node/.style={font=\scriptsize\sffamily}]
  \draw[black!55, line width=0.7pt] (0,-0.3) -- (0,5.4);
  \draw[black!55, line width=0.7pt] (0,-0.3) -- (9.2,-0.3);
  \node[anchor=east, text=black!60, rotate=90, anchor=south] at (-0.55,2.4) {Rate level};
  \node[anchor=north, text=black!60] at (4.6,-0.45) {Time};

  \draw[FRMGold, dashed, line width=1.2pt] (0.2,4.9) -- (6.6,5.15);
  \node[anchor=west, text=FRMGold] at (6.8,5.15) {$\mu$ \ \textit{(constant long-run target)}};

  \draw[FRMNavy, dashed, line width=1.2pt] (0.2,3.5) -- (6.6,3.85);
  \node[anchor=west, text=FRMNavy] at (6.8,3.85)
    {\textbf{LT} \ \textit{inflation, long-term trends}};

  \draw[FRMGreen, dashed, line width=1.2pt] (0.2,2.1) -- (6.6,2.5);
  \node[anchor=west, text=FRMGreen] at (6.8,2.5)
    {\textbf{MT} \ \textit{business cycles, monetary policy}};

  \draw[FRMRed, dashed, line width=1.2pt] (0.2,0.6) -- (6.6,0.85);
  \node[anchor=west, text=FRMRed] at (6.8,0.85) {\textbf{ST} \ \textit{fed funds rate}};
\end{tikzpicture}
\end{center}
\figcap{Four levels, each pulled toward the one above it. The short-term rate reacts quickly to
the medium-term factor, the medium-term factor adjusts more gradually to the long-term factor,
and the long-term factor shifts slowest of all toward the constant $\mu$.}

\begin{mrkeybox}[title={Note on the medium-term factor}]
Medium-term rates are affected in \term{two} ways:
\begin{enumerate}
\item Effects from short-term rates.
\item Their own independent volatility.
\end{enumerate}
\end{mrkeybox}

% ---------------------------------------------------------
\lo{MR-16 a}{Describe the structure of the Gauss+ model and discuss the implications of this structure for the model's ability to replicate empirically observed interest rate dynamics.}

Compared to the Vasicek model, the Gauss+ model is a more sophisticated \term{multifactor}
framework that provides a better representation of real-world interest rate behaviours.

\begin{mrfmlbox}
\begin{align*}
dr_t &= -\alpha_r\,(r_t - m_t)\,dt\\[3pt]
dm_t &= -\alpha_m\,(m_t - l_t)\,dt \;+\; \sigma_m\!\left(\rho\,dW_t^{1}
        + \sqrt{1-\rho^{2}}\;dW_t^{2}\right)\\[3pt]
dl_t &= -\alpha_l\,(l_t - \mu)\,dt \;+\; \sigma_l\,dW_t^{1}
\end{align*}
\mrvk{$r_t$, $m_t$, $l_t$ = the short-, medium- and long-term factors;\quad $\mu$ = the constant
long-run target;\quad $\alpha_r$, $\alpha_m$, $\alpha_l$ = the three mean reversion speeds;\quad
$\sigma_m$, $\sigma_l$ = volatilities of the medium- and long-term factors;\quad $\rho$ = the
correlation between them.}{%
\par\smallskip In the $dm_t$ equation, $\rho\,dW_t^{1}$ is the randomness \textbf{associated
with} the long-term factor, and $\sqrt{1-\rho^{2}}\,dW_t^{2}$ is the portion of randomness that
is \textbf{independent} of it.}
\end{mrfmlbox}

\begin{mrtrapbox}[title={The single most examinable feature}]
There is \term{no volatility term} in the $dr_t$ equation. The short-term rate has \textit{no
random shock of its own}; it only drifts toward $m_t$. That is deliberate: it reflects a central
bank pegging the policy rate, and it is what keeps short-rate volatility low so the model can
produce a \textbf{hump-shaped} volatility term structure instead of a declining one.

\medskip
It is also where the name comes from. ``Gauss'' means rates are normally distributed. The
``$+$'' flags that the model has \textbf{three factors but only two sources of risk}, $dW^1$ and
$dW^2$ --- an unusual state variable that is not itself a source of risk.
\end{mrtrapbox}

\subsection{The cascade of convergence}

\begin{center}
\begin{tikzpicture}[every node/.style={font=\scriptsize\sffamily},
  fac/.style={rectangle, rounded corners=3pt, line width=0.9pt,
    text width=2.55cm, align=center, inner sep=6pt, minimum height=1.25cm}]
  \node[fac, draw=FRMRed!60, fill=PaleRed] (r) at (0,0)
    {\textbf{$r_t$} short-term\\ \textit{fed funds}};
  \node[fac, draw=FRMGreen!60, fill=PaleGreen] (m) at (3.7,0)
    {\textbf{$m_t$} medium-term\\ \textit{business cycle}};
  \node[fac, draw=FRMNavy!60, fill=PaleNavy] (l) at (7.4,0)
    {\textbf{$l_t$} long-term\\ \textit{inflation}};
  \node[fac, draw=FRMGold!70, fill=PaleGold] (mu) at (11.1,0)
    {\textbf{$\mu$} constant\\ \textit{long-run target}};
  \draw[FRMGold, -{Stealth[length=5pt]}, line width=1.1pt] (r) -- (m);
  \draw[FRMGold, -{Stealth[length=5pt]}, line width=1.1pt] (m) -- (l);
  \draw[FRMGold, -{Stealth[length=5pt]}, line width=1.1pt] (l) -- (mu);
  \node[anchor=north, text=black!65, align=center] at (1.85,-0.75)
    {pulled toward, \textbf{fastest}\\ half-life 0.66 yr};
  \node[anchor=north, text=black!65, align=center] at (5.55,-0.75)
    {pulled toward, \textbf{slower}\\ half-life 1.09 yr};
  \node[anchor=north, text=black!65, align=center] at (9.25,-0.75)
    {pulled toward, \textbf{slowest}\\ half-life 42.0 yr};
\end{tikzpicture}
\end{center}
\figcap{Each factor is pulled toward the one to its right, and only $\mu$ is constant. If $r$ is
below $m$ it rises; if $r$ is above $m$ it falls; if $r = m$ there is no change in the short
rate. The half-lives are from the estimated parameters in MR-16 d and show how far apart the
three speeds are: under a year, about a year, and four decades.}

\begin{mrkeybox}[title={Key model characteristics}]
\begin{itemize}
\item \term{Mean reversion:} a cascade of convergence, $r \rightarrow m \rightarrow l
      \rightarrow \mu$. Only $\mu$ is constant; the other factors are dynamic.
\item \term{Speed of adjustment:} short-term rates react fastest, medium-term factors adjust
      more gradually, long-term factors shift slowest.
\item \term{Correlation:} changes in $m$ and $l$ are linked through the parameter $\rho$.
\item \term{Volatility:} the model can be calibrated to match the \textbf{hump-shaped}
      volatility term structure observed in markets --- low short-term volatility, higher
      medium-term volatility, declining long-term volatility.
\item \term{Pricing:} the model can be used to price bonds, swaps, caps, and floors.
\end{itemize}
\end{mrkeybox}

\begin{center}
\begin{tikzpicture}
\begin{axis}[
  width=10.8cm, height=5.4cm,
  xmin=0, xmax=31, ymin=0, ymax=90,
  xlabel={\scriptsize Term (years)},
  ylabel={\scriptsize Yield volatility (basis points)},
  xlabel style={font=\scriptsize}, ylabel style={font=\scriptsize},
  tick label style={font=\scriptsize},
  xtick={0,5,10,15,20,25,30}, ytick={0,20,40,60,80},
  legend style={font=\scriptsize\sffamily, draw=FRMRule, fill=white,
    at={(0.5,-0.30)}, anchor=north, legend columns=-1, column sep=10pt},
  axis line style={draw=black!55}, clip=false]
\addplot[FRMRed, line width=1.3pt, dashed, smooth, domain=0.5:30, samples=100]
  {85*exp(-0.075*x)};
\addlegendentry{Vasicek --- declining}
\addplot[FRMNavy, line width=1.3pt, smooth] coordinates {
(0.5,28)(1,38)(2,54)(3,64)(4,70)(5,73)(6,74)(7,73.5)(8,72)(10,68)
(12,64)(15,58)(20,50)(25,44)(30,39)};
\addlegendentry{Gauss+ --- hump-shaped, as observed}
\end{axis}
\end{tikzpicture}
\end{center}
\figcap{Why the missing volatility term matters. A mean-reverting single-factor model like
Vasicek can only produce a volatility curve that falls from its highest point at the short end.
Empirical and implied volatility curves hump: low at the very short end, peaking at intermediate
terms, then declining. Suppressing the short-rate shock is what lets Gauss+ reproduce that
shape.}

% ---------------------------------------------------------
\lo{MR-16 b}{Compare and contrast the dynamics, features, and applications of the Vasicek model and the Gauss+ model.}

\begin{center}
\small
\begin{tabularx}{\linewidth}{@{}p{2.6cm} >{\raggedright\arraybackslash}X >{\raggedright\arraybackslash}X@{}}
\arrayrulecolor{FRMRule}\toprule
\rowcolor{FRMNavy} \hdr{} & \hdr{Vasicek} & \hdr{Gauss+}\\
\textbf{Factors} & One: the short-term rate. & Three: $r$, $m$, $l$ --- but only \textbf{two}
  sources of risk.\\
\rowcolor{FRMSoft}
\textbf{Short-rate shock} & Present: $\sigma\,dw$. & \textbf{Absent.} $r$ only drifts toward
  $m$.\\
\textbf{Volatility term structure} & Declining, highest at the short end. &
  \textbf{Hump-shaped}, matching what is observed.\\
\rowcolor{FRMSoft}
\textbf{Dynamics} & Single mean reversion to a constant $\theta$. & A cascade: $r \to m \to l
  \to \mu$, at three very different speeds.\\
\textbf{Suited to} & Simple applications: pricing zero-coupon bonds, basic hedging. &
  Advanced applications: yield curve modeling, derivative pricing, risk management.\\
\rowcolor{FRMSoft}
\textbf{Weakness} & Cannot replicate complex term structure shapes, observed volatility across
  maturities, or macroeconomic and policy shifts. & More parameters to estimate; more complex to
  implement.\\
\bottomrule
\end{tabularx}
\end{center}

\begin{mrkeybox}
The Vasicek model is a \term{simpler framework} suitable for basic applications but suffers from
significant limitations in modeling real-world interest rate dynamics, particularly for
derivative pricing and term structure flexibility. The Gauss+ model addresses these deficiencies
with its multi-factor structure, cascading dynamics, and enhanced ability to fit market data.
\end{mrkeybox}

% ---------------------------------------------------------
\lo{MR-16 c}{Calculate changes in the short-term, medium-term, and long-term interest rate factors under the Gauss+ model.}

\noindent\gapbadge\hspace{4pt}{\sffamily\fontsize{7.6}{9}\selectfont\color{black!55}%
The source notes jump from objective 16.b straight to 16.d. This objective is absent from both
layers and is written below from the GARP curriculum.}\par
\vspace{6pt}

\begin{mrgapbox}[title={What the objective asks for}]
Given today's three factors, the estimated parameters, and a realization of the two random
shocks, compute $dr$, $dm$ and $dl$ over a short interval. The three equations from MR-16 a are
the whole of it; the work is bookkeeping, and the trap is which shock enters which equation.
\end{mrgapbox}

\begin{mrkeybox}[title={Read each equation before plugging in}]
\begin{itemize}
\item $dr$ has \term{no random term at all}. It is pure drift.
\item $dl$ is driven by $dW^{1}$ \term{only}.
\item $dm$ is driven by \term{both} shocks, weighted $\rho$ and $\sqrt{1-\rho^{2}}$, which is
      what gives it correlation $\rho$ with the long-term factor while keeping its own total
      volatility equal to $\sigma_m$.
\end{itemize}
\end{mrkeybox}

\begin{mrexbox}[title={Example --- one month of factor changes}]
Today $r = 2.00\%$, $m = 3.00\%$, $l = 4.00\%$. Parameters: $\alpha_r = 1.0547$,
$\alpha_m = 0.6358$, $\alpha_l = 0.0165$, $\sigma_m = 109.2$ bp, $\sigma_l = 96.4$ bp,
$\rho = 0.212$, $\mu = 10.555\%$. Over $dt = 1/12$, the shocks realize as $dW^{1} = 0.10$ and
$dW^{2} = -0.05$.
\tcblower
\small
\textbf{Short-term factor} --- drift only:
\[ dr = -1.0547\,(2.00\% - 3.00\%)\left(\tfrac{1}{12}\right) = \mathbf{+0.0879\%} \]
The short rate is \textit{below} the medium-term factor, so it is pulled \textbf{up}.

\medskip
\textbf{Medium-term factor} --- drift plus both shocks, with
$\sqrt{1 - 0.212^{2}} = 0.9773$:
\begin{align*}
\text{drift} &= -0.6358\,(3.00\% - 4.00\%)\left(\tfrac{1}{12}\right) = +0.0530\%\\
\text{shock} &= 1.092\%\bigl[(0.212)(0.10) + (0.9773)(-0.05)\bigr] = -0.0302\%\\
dm &= 0.0530\% - 0.0302\% = \mathbf{+0.0228\%}
\end{align*}
The drift pulls it up toward $l$, but the shock partly offsets.

\medskip
\textbf{Long-term factor} --- drift plus $dW^{1}$ only:
\begin{align*}
\text{drift} &= -0.0165\,(4.00\% - 10.555\%)\left(\tfrac{1}{12}\right) = +0.0090\%\\
\text{shock} &= 0.964\% \times 0.10 = +0.0964\%\\
dl &= 0.0090\% + 0.0964\% = \mathbf{+0.1054\%}
\end{align*}
\end{mrexbox}

\begin{mrtrapbox}[title={What the three answers show}]
The long-term factor moves the \textit{most} this month, at 10.5 basis points, despite having by
far the slowest mean reversion. That is because almost all of its move is the random shock ---
its drift contributes under one basis point, since $\alpha_l$ is tiny. Meanwhile the short rate
moves 8.8 basis points with no randomness at all. \term{Slow mean reversion does not mean a
small move}; it means the factor is not pulled back quickly once it moves.
\end{mrtrapbox}

% ---------------------------------------------------------
\lo{MR-16 d}{Explain how the parameters of the Gauss+ model can be estimated from empirical data.}

The parameters can be estimated using \term{maximum likelihood} techniques. However, there is a
simplified approach that employs the following steps. This procedure is possible when random
components are not used to determine mean reversion speeds.

\begin{center}
\small
\begin{tabularx}{\linewidth}{@{}p{1.4cm} >{\raggedright\arraybackslash}X@{}}
\arrayrulecolor{FRMRule}\toprule
\rowcolor{FRMNavy} \hdr{Step} & \hdr{What you do}\\
\textbf{1} & Select a \term{proxy for the short-term rate} --- the fed funds rate, for
  example.\\
\rowcolor{FRMSoft}
\textbf{2} & \term{Regress} the changes in the short-term rate against changes in a medium-term
  rate (a 2-year rate) and a longer-term rate (a 10-year rate).\\
\textbf{3} & Find the mean reversion speeds $\alpha_r$, $\alpha_m$ and $\alpha_l$ that replicate
  the sensitivities seen in Step 2. This is possible because the regression coefficients
  \term{do not depend on the volatility parameters}, only on the mean reversion parameters.\\
\rowcolor{FRMSoft}
\textbf{4} & From market data, determine the \term{volatility term structure}. This can be found
  from the observed volatility of rates, or from the volatilities the market uses to price
  interest rate derivatives.\\
\textbf{5} & Find the volatility and correlation parameters $\sigma_m$, $\sigma_l$ and $\rho$
  such that the volatility term structure from Step 4 is replicated.\\
\rowcolor{FRMSoft}
\textbf{6} & Find the long-term reversion level $\mu$ by \term{minimizing the squared errors} of
  observed yields versus estimated yields.\\
\bottomrule
\end{tabularx}
\end{center}

\begin{mrkeybox}[title={Why the staging works}]
The estimation splits cleanly because the model's regression coefficients depend only on the
mean reversion parameters and not on the volatility parameters. That lets Steps 2 and 3 pin down
the $\alpha$'s before any volatility is touched, and Steps 4 and 5 then pin down the
$\sigma$'s and $\rho$ with the $\alpha$'s already fixed.
\end{mrkeybox}

\subsection{The estimated parameters, for reference}

\begin{center}
\small
\begin{tabular}{l r r}
\arrayrulecolor{FRMRule}\toprule
\rowcolor{FRMNavy} \hdr{Parameter} & \hdr{Estimate} & \hdr{Half-life (years)}\\
$\alpha_r$ & 1.0547 & 0.66\\
\rowcolor{FRMSoft} $\alpha_m$ & 0.6358 & 1.09\\
$\alpha_l$ & 0.0165 & 42.01\\
\midrule
\rowcolor{FRMSoft} $\sigma_m$ & 109.2 bp & \\
$\sigma_l$ & 96.4 bp & \\
\rowcolor{FRMSoft} $\rho$ & 0.212 & \\
$\mu$ & 10.555\% & \\
\bottomrule
\end{tabular}
\end{center}
\figcap{Estimated from US Treasury zero coupon yields, January 2014 to January 2022. The mean
reversion speeds come out in exactly the order the economics predicts: the central bank reaction
is fastest, the medium-term factor's reversion next, and the long-term factor's reversion to
$\mu$ slowest by a wide margin.}

\begin{center}
\begin{tikzpicture}
\begin{axis}[
  width=10.4cm, height=4.4cm,
  xmin=0.15, xmax=100, xmode=log, ymin=0, ymax=4,
  xlabel={\scriptsize Half-life in years (log scale)},
  axis y line=none, axis x line=bottom, axis line style={draw=black!55},
  xlabel style={font=\scriptsize}, tick label style={font=\scriptsize},
  xtick={0.25,0.5,1,2,5,10,25,50,100},
  xticklabels={0.25,0.5,1,2,5,10,25,50,100},
  clip=false]
\addplot[only marks, mark=*, mark size=2.6pt, color=FRMRed] coordinates {(0.66,3)};
\addplot[only marks, mark=*, mark size=2.6pt, color=FRMGreen] coordinates {(1.09,2)};
\addplot[only marks, mark=*, mark size=2.6pt, color=FRMNavy] coordinates {(42.01,1)};
\node[anchor=west, font=\scriptsize\sffamily, text=FRMRed] at (axis cs:0.78,3)
  {$\alpha_r$ --- short rate, \textbf{0.66 yr}};
\node[anchor=west, font=\scriptsize\sffamily, text=FRMGreen] at (axis cs:1.30,2)
  {$\alpha_m$ --- medium-term factor, \textbf{1.09 yr}};
\node[anchor=east, font=\scriptsize\sffamily, text=FRMNavy] at (axis cs:35,1)
  {$\alpha_l$ --- long-term factor, \textbf{42.0 yr}};
\end{axis}
\end{tikzpicture}
\end{center}
\figcap{The three half-lives on a log scale, because on a linear scale the first two would sit
on top of each other. The long-term factor reverts roughly sixty times more slowly than the
short rate, which is the quantitative content of the phrase ``cascade of convergence''.}


% ==========================================================
